\documentclass[11pt]{article}
% =========================================================
%  學生講義 共用前言（以 XeLaTeX 編譯）
%  需要套件：xeCJK, tcolorbox（其餘多在 BasicTeX 內）
% =========================================================
\usepackage[a4paper,margin=2.3cm]{geometry}
\usepackage{amsmath,amssymb}
\usepackage{xcolor}
\usepackage{graphicx}

% ---- 中文字型（macOS 系統字型）----
\usepackage{xeCJK}
\setCJKmainfont{Songti TC}      % 內文：宋體
\setCJKsansfont{Heiti TC}       % 標題/強調：黑體
\xeCJKsetup{AutoFakeBold=2.2, AutoFakeSlant=0.2}

% ---- 版面 ----
\setlength{\parindent}{0pt}
\setlength{\parskip}{0.55em}
\linespread{1.15}
\renewcommand{\baselinestretch}{1.15}

% ---- 顏色 ----
\definecolor{cBlue}{HTML}{1f6feb}
\definecolor{cGray}{HTML}{6b7280}
\definecolor{cGrayBg}{HTML}{f3f4f6}
\definecolor{cPurp}{HTML}{7a3ff2}
\definecolor{cPurpBg}{HTML}{f5f1ff}

% ---- 標註框 ----
\usepackage[most]{tcolorbox}

% 就地補充新概念（灰底）：\begin{補充框}{標題}
\newtcolorbox{補充框}[1]{
  enhanced, breakable, arc=2pt, boxrule=0.6pt,
  colback=cGrayBg, colframe=cGray,
  left=9pt,right=9pt,top=7pt,bottom=7pt,
  before upper={\textbf{\sffamily #1}\par\vspace{3pt}},
}

% 延伸 / 開放問題（紫色左邊條）：\begin{延伸框}{標題}
\newtcolorbox{延伸框}[1]{
  enhanced, breakable, arc=2pt, boxrule=0pt,
  colback=cPurpBg, colframe=cPurpBg,
  borderline west={3pt}{0pt}{cPurp},
  left=10pt,right=9pt,top=7pt,bottom=7pt,
  before upper={\textbf{\sffamily 延伸閱讀｜#1}\par\vspace{3pt}},
}

% ---- 圖：置中、底下小字說明 ----
% \fig{檔名}{寬度}{說明}
\newcommand{\fig}[3]{%
  \begin{center}
  \includegraphics[width=#2\linewidth]{#1}\\[2pt]
  {\small\color{cGray} #3}
  \end{center}}

% ---- 章節標題用黑體 ----
\newcommand{\h}[1]{\sffamily #1}


\begin{document}

\begin{center}
{\sffamily\Huge\bfseries 必物-5 能量}\\[3pt]
{\sffamily\large 普通高中 物理（必修）・學生講義}
\end{center}
\vspace{0.6em}

能量（energy）是貫穿整個物理的共同量。本章先認識能量的各種形式及其相互轉換，並掌握最重要的規律——\textbf{總能量守恆}；接著進入二十世紀的重要發現：質量本身也是能量（$E=mc^2$），由此理解太陽與核能的來源；最後回到日常的熱，釐清\textbf{溫度、熱、內能}三者的分別，以及一條決定「時間方向」的規律——熱不能完全變回功。

本章以\textbf{定性與半定量}理解為主，完整的功能定量計算（功、動能定理、力學能守恆）留待高二。本章談的微觀能量與量子化伏筆，直通\textbf{必物-6 量子現象}。

\medskip
本章主線：

\begin{center}
能量的形式與守恆 $\rightarrow$ 質能等價 $E=mc^2$ $\rightarrow$ 核能（融合與分裂）$\rightarrow$ 溫度、熱與內能 $\rightarrow$ 熱的方向性（第二定律）
\end{center}

\section*{\sffamily 5-1\quad 能量的形式、轉換與守恆}

\subsection*{\sffamily A.\ 能量的形式}

能量不易用一句話定義，但可用一個直觀說法：\textbf{能量是讓事情發生（讓物體運動、加熱、發光、起化學反應）的本錢}。一個系統能量越多，能造成的改變就越多。物理把能量分成幾種常見的\textbf{形式}：

\begin{center}
\renewcommand{\arraystretch}{1.3}
\begin{tabular}{p{0.26\linewidth} p{0.30\linewidth} p{0.30\linewidth}}
\hline
\textbf{形式} & \textbf{存在於} & \textbf{生活例子} \\
\hline
動能（kinetic energy） & 運動的物體 & 行駛的車、流動的水 \\
位能（potential energy） & 物體在力場中的位置、或形變 & 高處的水、拉長的彈簧 \\
熱能（thermal energy） & 分子的雜亂運動 & 燒燙的水、摩擦生熱 \\
化學能（chemical energy） & 原子間的鍵結 & 食物、汽油、電池 \\
電能／電磁能（electromagnetic energy） & 電場與磁場、電流、光 & 插座電力、陽光、訊號 \\
核能（nuclear energy） & 原子核內部 & 太陽、核電廠 \\
\hline
\end{tabular}
\end{center}

\begin{補充框}{能量守恆（conservation of energy）}
能量\textbf{不會無中生有，也不會憑空消失}，只會\textbf{從一種形式轉換成另一種形式}；一個孤立系統內，各種形式能量的\textbf{總和保持不變}：
$$\boxed{\;E_{\text{總}} = \sum_i E_i = \text{定值}\;}$$
這是物理學中適用範圍最廣的定律之一，從力學、電磁到量子皆成立。

\textbf{注意：}「能量被『用掉』了」並不精確。能量從不消失，只是轉成比較不易再利用的形式（多半是散逸到環境的低溫熱能）。所謂「能源危機」，精確說法是「\textbf{可用的、有品質的能量}越來越少」，而非能量總量減少。這個「品質」概念在 5-5 會正式登場。
\end{補充框}

幾乎所有日常裝置，本質上都是\textbf{能量轉換器}：

\begin{itemize}\setlength{\itemsep}{1pt}
\item \textbf{水力發電}：水庫的位能 $\rightarrow$ 落水的動能 $\rightarrow$ 發電機把動能變電能 $\rightarrow$ 燈泡把電能變光與熱。
\item \textbf{燃燒汽油}：汽油的化學能 $\rightarrow$ 引擎的熱能 $\rightarrow$ 活塞與車輪的動能。
\item \textbf{植物}：太陽的光能 $\rightarrow$ 光合作用儲成化學能。
\end{itemize}

每一步都有能量轉換，但\textbf{全程加總，總能量不變}。

\fig{m5-conversion}{0.62}{圖 5-1：能量在各種形式之間轉換，但總量始終守恆。}

\textbf{注意：}「能量守恆」不等於「力學能守恆」。力學能（動能＋位能）守恆只在\textbf{沒有摩擦等耗散}時成立；一旦有摩擦，力學能會減少，但它變成了熱能，\textbf{總能量}仍守恆。

\subsection*{\sffamily B.\ 同一份能量，在不同尺度有不同面貌}

能量「形式」的劃分，其實取決於\textbf{從哪個尺度觀察}。一杯熱水，從整杯的\textbf{巨觀}尺度看靜止不動、沒有動能；但放大到\textbf{微觀}尺度，每個水分子都在劇烈運動——所謂「熱能」，在微觀正是\textbf{大量分子的動能與分子間交互作用能的總和}。同樣地，化學能在微觀上是電子與原子核之間的電磁位能（鍵結能），彈簧位能在微觀上是原子被擠壓或拉開後的電磁位能。

換個尺度看，熱能就是動能、化學能就是電磁能。能量並非真的分成六種不同「物質」。這個觀念使 5-4 的「熱能其實是微觀動能」變得自然。

\subsection*{\sffamily C.\ 電場、磁場也帶有能量；電磁波會搬運能量}

能量不一定要附在有形物體上。\textbf{只要空間中有電場或磁場，就帶有能量}：把電容器充電、或讓電流通過線圈，能量就儲存在電場、磁場裡。當電場與磁場一起振盪並向外傳播，便形成\textbf{電磁波（electromagnetic wave）}——它能在真空中傳播，並把能量從一處搬到另一處。例如太陽光就是電磁波，把太陽核心的能量送到地球；手機通訊則是把訊息載在電磁波上發射、由基地台接收。電磁波的產生機制（變動的電場生磁場、變動的磁場生電場）見必物-4。

\section*{\sffamily 5-2\quad 質量也是能量：$E=mc^2$}

\subsection*{\sffamily A.\ 質能等價}

在 1905 年之前，\textbf{質量守恆}與\textbf{能量守恆}被當成兩條互不相干的獨立定律。愛因斯坦（Einstein）的狹義相對論揭示：質量與能量是同一件事的兩種說法，可以互相轉換，換算係數是光速的平方。

\begin{補充框}{$E=mc^2$ 的意義：質量本身就是能量的一種形式}
一個質量為 $m$ 的物體，本身就蘊含能量
$$\boxed{\;E = mc^2\;}\qquad(c=3\times10^8\ \text{m/s}\ \text{為光速})$$
這稱為物體的\textbf{靜止能量（rest energy）}。

$E=mc^2$ 真正的意思\textbf{不是}「質量變成能量」，而是：質量本身就是能量的一種形式。一個靜止物體，光是存在就帶有能量 $E=mc^2$。質量與能量並非兩種東西再去互相兌換，而是\textbf{同一件事}；$c^2$ 只是因為質量用 kg、能量用 J，兩種單位之間的換算係數。可類比為「1 美元 = 30 元台幣」：美元和台幣是同一筆錢的兩種計法，不是美元「變成」台幣。

\textbf{所謂「質量虧損」是什麼？}以核融合為例：氘＋氚 $\rightarrow$ 氦＋中子。若很精確地測量，會發現
$$(\text{反應後產物的總質量})\;<\;(\text{反應前的總質量}).$$
差的這一點點質量 $\Delta m$ 並沒有消失，而是以\textbf{能量}的形式（產物的動能、放出的光等）釋出，數值剛好是 $E=\Delta m\,c^2$。換個角度看：\textbf{綁得越緊的系統，總質量越小}。把核子綁進原子核會放出能量，所以「綁好的核」比「散開的核子」輕，輕的量等於放出的能量除以 $c^2$。可類比為把幾顆磁鐵從遠處吸在一起：吸合時放出能量（撞擊聲、發熱），組好的磁鐵組合能量比分開時低。
\end{補充框}

\subsection*{\sffamily B.\ 為什麼「一點點質量」對應「巨大能量」}

關鍵在 $c^2$ 這個係數非常大：
$$c^2=(3\times10^8)^2=9\times10^{16}\ \text{m}^2/\text{s}^2.$$
所以即使只有 1 克（$0.001$ kg）的質量完全轉成能量，得到的能量是
$$E=0.001\times9\times10^{16}=9\times10^{13}\ \text{J},$$
相當於約 2 萬噸 TNT，是一顆原子彈的量級。這就是為什麼核反應（牽涉到質量的改變）能釋放出化學反應（只重排電子、質量幾乎不變）遠遠不及的能量。

\textbf{注意：}$E=mc^2$ 並非只在核反應才適用。\textbf{任何}放出能量的過程，質量都會減少，只是極微小。拉開一個彈簧、燒一根火柴、電池放電，產物的總質量都比反應物少了一點點（少的量＝放出的能量 $\div c^2$），只是小到秤不出來。質能等價是\textbf{普適}的；過去以為「化學反應質量守恆」，是因為這個變化太小，任何天平都測不出。

\section*{\sffamily 5-3\quad 核能：融合與分裂}

\subsection*{\sffamily A.\ 束縛能}

原子核由\textbf{質子與中子}（合稱\textbf{核子}）組成。要把一個原子核拆散成一顆顆自由核子，必須輸入能量；反過來，把自由核子組成原子核時，會放出同樣多的能量。這份能量稱為\textbf{束縛能（binding energy）}。由質能等價，放出能量代表系統變輕，因此「綁好的核」比「散開的核子」輕——這正是 5-2 所述質量虧損的來源。

重點不是「整個核的束縛能」，而是\textbf{平均每個核子分到多少束縛能}。把它對質量數 $A$（核子數）作圖，得到核物理中重要的一張圖：束縛能曲線。

\begin{補充框}{每核子束縛能與束縛能曲線}
\textbf{每核子束縛能}＝整個核的束縛能除以核子數 $A$，代表「平均每個核子被綁得多緊」。這個量越大，核越穩定。曲線的關鍵：
\begin{itemize}\setlength{\itemsep}{1pt}
\item 曲線在\textbf{鐵（$^{56}$Fe）附近達到最高峰}（每核子束縛能約 8.8 MeV），代表鐵是\textbf{最穩定}的核。
\item \textbf{峰的左邊（輕核）}：把兩個輕核\textbf{融合}成較重的核，會更靠近峰頂 $\rightarrow$ 放出能量。
\item \textbf{峰的右邊（重核）}：把一個重核\textbf{分裂}成兩個中等核，也會更靠近峰頂 $\rightarrow$ 放出能量。
\end{itemize}
因此\textbf{融合與分裂都能釋放能量}，只是分屬曲線兩側、往同一個「最穩定的鐵峰」靠攏。
\end{補充框}

\fig{m5-binding}{0.78}{圖 5-2：束縛能曲線。縱軸為每核子平均束縛能、橫軸為質量數 $A$。鐵（$^{56}$Fe）位於峰頂、最穩定；左側輕核融合、右側重核分裂，都往峰頂靠攏並釋放能量。}

\subsection*{\sffamily B.\ 核分裂（fission）}

一個\textbf{重核}（如鈾 $^{235}$U）吸收一個中子後分裂成兩個中等質量的核，同時放出 2$\sim$3 個中子與大量能量。放出的中子可再去打中下一個鈾核，形成\textbf{連鎖反應（chain reaction）}：

\begin{itemize}\setlength{\itemsep}{1pt}
\item \textbf{緩慢、受控}的連鎖反應 $\rightarrow$ \textbf{核能發電廠}（用控制棒吸收多餘中子，穩定輸出）。
\item \textbf{瞬間、失控}的連鎖反應 $\rightarrow$ 原子彈。
\end{itemize}

\subsection*{\sffamily C.\ 核融合（fusion）}

兩個\textbf{輕核}（如氫的同位素氘 $^2$H、氚 $^3$H）融合成較重的核（如氦 $^4$He），放出能量。\textbf{太陽與所有恆星}的能量來源正是核融合（氫融合成氦）。融合釋放的能量（按每核子計）比分裂更多，且原料便宜、產物乾淨（氦無放射性）。但融合\textbf{極難點燃}：兩個帶正電的核要靠得夠近才能融合，必須克服強烈的電磁排斥，需要上億度的高溫高壓。人類目前仍在研究如何讓「受控核融合」穩定地\textbf{淨產出}能量（如國際 ITER 計畫）。

\textbf{注意：}關於核能有幾點常見的混淆需要釐清——

\begin{itemize}\setlength{\itemsep}{1pt}
\item \textbf{融合與分裂兩者都釋放能量}，差別只在原料位於束縛能曲線的哪一側；都往鐵峰靠攏就放能。並非一者釋能、一者吸能。
\item \textbf{太陽不是化學燃燒}。化學燃燒需要氧氣，且能量遠遠不夠；太陽靠的是核融合。
\item \textbf{核電廠不會像原子彈一樣核爆}。核電廠燃料的鈾濃度遠低於核彈，物理上不可能發生核爆；其風險是爐心過熱熔毀與放射性外洩，與核爆是兩回事。
\end{itemize}

核能是能量密度極高、發電過程不排碳，但具有放射性廢料長期處置、重大事故風險與核武擴散疑慮的能源議題。

\section*{\sffamily 5-4\quad 溫度、熱與內能}

\subsection*{\sffamily A.\ 內能}

把觀察尺度推進到物體內部：原子、分子\textbf{從不靜止}，它們不停地振動、平移、轉動（\textbf{微觀動能}），彼此之間又有鍵結與排斥（\textbf{微觀位能／交互作用能}）。

\begin{補充框}{內能（internal energy）與熱能}
物體內部所有分子的\textbf{微觀動能}與分子間\textbf{交互作用位能}的總和，稱為\textbf{內能}；課綱中把這份「原子的動能與交互作用能的總和」稱為\textbf{熱能}。它是物體\textbf{內部本身擁有}的能量，與物體整體的運動、所在高度無關。
\end{補充框}

\subsection*{\sffamily B.\ 溫度}

\begin{補充框}{溫度（temperature）：分子平均動能的量度}
\textbf{溫度}是分子\textbf{平均微觀動能}的量度。溫度越高，分子平均動得越劇烈。溫度是一個\textbf{強度量}，與分子有多少個無關。對\textbf{理想氣體}，分子間幾乎無交互作用，內能幾乎全是分子動能，因此
$$E_{\text{內}}\;\propto\;T\quad(\text{理想氣體}).$$
此處 $T$ 必須是\textbf{絕對溫度}（見下方克氏溫標），不能代入攝氏度數。
\end{補充框}

\fig{m5-internal}{0.78}{圖 5-3：溫度是分子亂動劇烈程度的量度。左側低溫、分子慢、動能小；右側高溫、分子快、動能大。對理想氣體，內能 $\propto T$。}

\subsection*{\sffamily C.\ 克氏溫標與絕對零度}

日常用的攝氏溫標把水的冰點訂為 $0\,^\circ$C，是人為的方便基準。物理上更自然的是\textbf{克氏溫標（Kelvin scale，絕對溫標）}，它把「分子幾乎完全停止亂動、微觀動能降到最低」的狀態訂為起點：
$$\boxed{\;T(\text{K})=t(^\circ\text{C})+273.15\;}$$
這個起點 $0\ \text{K}\approx-273.15\,^\circ$C 稱為\textbf{絕對零度（absolute zero）}，是溫度的\textbf{理論下限}。克氏溫標沒有負值。

\begin{補充框}{絕對零度為什麼是溫度的下限？能不能更冷？}
溫度量度的是\textbf{分子平均微觀動能}。動能 $\tfrac12mv^2$ 不可能是負的，最小就是\textbf{零}（分子完全不動）。既然平均動能有個最小值（零），對應的溫度自然也有個最小值，那就是絕對零度。

往高溫的方向，分子可以動得越來越快、動能越來越大，\textbf{沒有上限}；往低溫的方向，最多就是「不動」，再也減不下去——所以\textbf{溫度有下限、沒有上限}。把絕對零度設為刻度的 0，克氏溫度便直接正比於分子平均動能，這也是「理想氣體內能 $\propto T$」必須用克氏溫標的原因。

\textbf{注意：}有兩點容易誤解。第一，嚴格說「絕對零度時分子完全靜止」並不精確——量子力學指出粒子即使在最低能量狀態仍保有一點\textbf{零點能（zero-point energy）}的抖動，絕對零度對應的是「降到能量最低的量子態」，而非字面上的速度為零。第二，0 K 可以\textbf{無限逼近、卻永遠達不到}：熱力學第三定律指出，要冷卻到恰好 0 K 需要無窮多個步驟。實驗室已能冷到十億分之幾 K（比星際空間還冷），仍差那麼一點碰不到 0。
\end{補充框}

\subsection*{\sffamily D.\ 熱}

最容易混淆的一點是：\textbf{「熱」和「內能」不是同一回事}。

\begin{補充框}{熱（heat）：傳遞中的能量，是過程量}
\textbf{由於兩物體溫度不同而從高溫物體流向低溫物體的能量，稱為熱}。熱是一個\textbf{過程量}——它描述能量\textbf{正在傳遞}這件事，只有在能量流動時才談得上「熱」。物體本身\textbf{擁有}的是\textbf{內能}，不是「熱」。

可類比為金錢：物體擁有的是\textbf{內能}（像戶頭裡的存款），而\textbf{熱}是\textbf{轉帳的動作}（錢正從一個戶頭流到另一個戶頭）。轉完帳，「轉帳」就不再存在——它從來不是戶頭裡的東西，而是「流動」這件事本身。\textbf{溫度}則像\textbf{水位高低}，決定能量往哪邊流：熱總是從溫度高的一方流向低的一方，與哪邊內能多無關。
\end{補充框}

\textbf{注意：}以下三點是熱學最常見的混淆，需要分清楚——

\begin{itemize}\setlength{\itemsep}{1pt}
\item \textbf{物體不「含有」熱}。物體含有的是\textbf{內能}；熱是溫差造成的能量\textbf{傳遞過程}，傳完就成為對方的內能，不再叫熱。
\item \textbf{溫度高不一定內能多}。一杯 $90\,^\circ$C 的水，內能可能比一整缸 $40\,^\circ$C 的洗澡水還少（因為洗澡水分子數多很多）。溫度看的是\textbf{平均}動能，內能看的是\textbf{總量}。
\item \textbf{沒有「冷」在流動}。是\textbf{熱（能量）從溫度高的一方流向低的一方}；覺得「冷氣灌進來」，其實是身上的熱被帶走。
\end{itemize}

熱的\textbf{流向只看溫度，不看內能}。把一杯 $90\,^\circ$C 的水倒進一缸 $40\,^\circ$C 的水，儘管大缸的內能遠大於小杯，熱仍從小杯（高溫）流向大缸（低溫），直到同溫為止。

\subsection*{\sffamily E.\ 熱平衡}

熱既然是因溫差而傳遞，就要追問：它什麼時候停？把一杯熱水和一杯冷水接觸，熱會從熱水流向冷水；但溫差越來越小，傳遞越來越慢，\textbf{直到兩者溫度相等}——這時不再有淨能量流動。

\begin{補充框}{熱平衡（thermal equilibrium）與熱的流向}
兩物體接觸時，熱\textbf{只會自發地從高溫物體流向低溫物體}，使高溫者降溫、低溫者升溫；當兩者\textbf{溫度相等}時，淨熱流停止，系統達到\textbf{熱平衡}。有兩點要分清楚：
\begin{itemize}\setlength{\itemsep}{1pt}
\item 熱的流向只由\textbf{溫度}（高 $\rightarrow$ 低）決定，\textbf{與內能總量無關}——高溫的小物體仍會把熱傳給低溫的大物體，即使後者內能遠大。
\item 平衡時是「\textbf{淨}」熱流為零，並非微觀上不再交換，而是兩個方向的交換達到平衡。
\end{itemize}
「熱只從高溫流向低溫、不會自發逆流」這個方向性，正是 5-5 要深入的熱力學第二定律。
\end{補充框}

\subsection*{\sffamily F.\ 熱傳遞的三種方式}

熱是「怎麼」從高溫處傳到低溫處的？標準上分成三種途徑，三者常同時發生。

\begin{補充框}{傳導、對流、輻射}
\begin{itemize}\setlength{\itemsep}{2pt}
\item \textbf{傳導（conduction）}：靠\textbf{分子／自由電子的碰撞}把動能一個傳一個地接力過去，物質本身\textbf{不整體移動}。固體（尤其金屬，靠自由電子）主要靠傳導。例：鐵湯匙一端泡熱湯，另一端也會變燙。
\item \textbf{對流（convection）}：靠\textbf{流體（液體、氣體）本身的整體流動}把熱「載」走。受熱的流體膨脹變輕而上升、冷流體下沉，形成循環。例：燒開水、暖氣使整個房間變暖、海陸風。
\item \textbf{輻射（radiation）}：靠\textbf{電磁波}傳熱，\textbf{不需要任何介質}，可在真空中進行。任何有溫度的物體都會發出熱輻射。例：太陽的熱橫越真空到地球、靠近火堆正面就覺得燙、紅外線測溫槍。
\end{itemize}
只有\textbf{輻射}能穿越真空——這也呼應 5-1「電磁波能在真空中搬運能量」。
\end{補充框}

\textbf{注意：}兩點常見的混淆——
\begin{itemize}\setlength{\itemsep}{1pt}
\item \textbf{對流不會在固體裡發生}。對流一定要有流體的整體流動，固體分子被鎖在晶格上無法整體流動，故固體只能靠傳導（與輻射）。
\item \textbf{真空中並非完全無法傳熱}。傳導與對流確實不行（沒有介質），但輻射可以：保溫瓶夾層抽真空，正是用真空切斷傳導與對流，再鍍反光層減少輻射，才能保溫。
\end{itemize}

\section*{\sffamily 5-5\quad 熱的方向性：為什麼熱不能完全變回功}

\subsection*{\sffamily A.\ 一個不對稱的事實}

5-1 說能量守恆、各種形式可以互相轉換。但\textbf{熱}的轉換是\textbf{不對稱}的：

\begin{itemize}\setlength{\itemsep}{1pt}
\item \textbf{功 $\rightarrow$ 熱：可以 100\% 轉換}。摩擦、攪拌、煞車，機械功總能毫無保留地全變成熱。
\item \textbf{熱 $\rightarrow$ 功：無法 100\% 轉換}。任何把熱變成功的機器（\textbf{熱機}），都\textbf{一定有一部分熱被排到低溫處}，不可能把吸收的熱全部變成功。
\end{itemize}

這條「功能全變熱、熱不能全變功」的經驗事實，就是\textbf{熱力學第二定律}的一種講法（本章只談定性概念）。

\fig{m5-engine}{0.92}{圖 5-4：能量的方向性。左——功可以 100\% 變成熱；右——熱機從高溫源吸熱 $Q_H$、對外作功 $W$，但一定有一部分熱 $Q_C$ 排到低溫源，效率小於 100\%。}

\subsection*{\sffamily B.\ 熱機與效率}

\textbf{熱機（heat engine）}從高溫源吸熱 $Q_H$，對外作功 $W$，再把剩下的熱 $Q_C$ 排到低溫源。由能量守恆 $Q_H=W+Q_C$，效率定義為「得到的功」佔「付出的熱」的比例：
$$\boxed{\;\eta=\frac{W}{Q_H}=\frac{Q_H-Q_C}{Q_H}=1-\frac{Q_C}{Q_H}\;}$$
第二定律告訴我們 $Q_C$ \textbf{不可能為零}，所以 $\eta$ 永遠\textbf{小於 1（小於 100\%）}。汽車引擎效率約 25$\sim$30\%、火力發電廠約 40\%，其餘都當廢熱排掉了。

\textbf{注意：}熱機效率達不到 100\% \textbf{不是工程不夠好}。即使是零摩擦、零散熱的\textbf{理想}熱機，效率仍受 $1-\dfrac{T_C}{T_H}$（卡諾上限）限制，不可能達到 100\%。這是自然律，不是技術問題。

\subsection*{\sffamily C.\ 能量的「品質」與不可逆}

第二定律其實在說：\textbf{能量除了「量」，還有「品質」}。

\begin{itemize}\setlength{\itemsep}{1pt}
\item \textbf{高品質能量}（功、電能、集中的高溫熱）：可以轉成幾乎任何其他形式。
\item \textbf{低品質能量}（散逸到環境的低溫廢熱）：很難再轉回有用的功。
\end{itemize}

每一次自然發生的過程，能量\textbf{總量守恆}，但\textbf{品質一路下降}——有用的功不斷退化成散亂的低溫熱，而且這個過程是\textbf{不可逆（irreversible）}的：打翻的水不會自己回到杯子、熱不會自己從冷的流向熱的。物理用\textbf{熵（entropy）}這個量來量度這種「散亂、不可用」的程度：孤立系統的熵只增不減。

\begin{延伸框}{為什麼功能全變熱、熱卻不能全變回功？熵與時間之箭}
功變熱與熱變功的差別，本質是\textbf{有序程度}的差別。\textbf{功（與電能）是有序的能量}：例如一個飛輪轉動，大量原子朝同一方向整齊運動。\textbf{熱是無序的能量}：高溫物體裡，大量分子朝四面八方亂動。於是：
\begin{itemize}\setlength{\itemsep}{1pt}
\item \textbf{功 $\rightarrow$ 熱}：把整齊的運動打散成亂動，可類比為把排好的撲克牌往空中一撒——自然發生。
\item \textbf{熱 $\rightarrow$ 功}：要把亂動重新收攏成整齊一致的運動，可類比為要撒落一地的牌自己飛回排好——不可能自發完成。
\end{itemize}
熱機之所以能把一部分熱變功，是靠\textbf{溫差}：讓熱從高溫流向低溫，順手抽出一小部分變成有序的功；但剩下的熱非排到低溫源不可，這部分回收不了。

物理用\textbf{熵} $S$ 量度系統的無序程度。第二定律最一般的講法是：\textbf{孤立系統的熵只會增加或維持不變，永不減少}（$\Delta S_{\text{孤立}}\ge 0$）。功變熱是有序變無序、熵增加，自發合法；熱完全變功是無序變有序、熵減少，違反第二定律。因此熱機效率上限是「總熵不准減少」逼出來的硬底線，而非工程問題。（局部熵仍可減少：冰箱、冷氣、生命體都能讓局部變有序，但代價是讓環境熵增加更多，總熵仍增加。）

第二定律有個獨特地位：\textbf{它是少數能區分「過去」與「未來」的物理定律}。牛頓力學、電磁學、量子力學的基本方程式幾乎都是\textbf{時間對稱}的——把時間倒著放，方程式照樣成立，一段撞球碰撞的影片正放、倒放都合理。唯獨牽涉到熱與熵的過程能看出時間方向：熱水自己變涼（熵增，正常）；涼水自己變燙、把熱送回爐子（熵減，一看就是影片倒放）。這個「時間只朝熵增方向流」的現象，稱為\textbf{時間之箭（arrow of time）}。

\textbf{這部分仍是開放問題。}「熵會增加」本身已是\textbf{定論}，無數實驗與統計力學都支持它：波茲曼指出「無序」的狀態數遠多於「有序」的狀態數，系統幾乎必然從少數有序狀態演化到絕大多數的無序狀態，熵增其實是「機率壓倒性地大」。但\textbf{時間之箭的終極起源目前尚無定論}：熵會增加，前提是宇宙\textbf{過去}處於一個熵極低的特殊狀態，為什麼宇宙誕生時熵這麼低，把問題推到了宇宙學與宇宙初始條件，至今並無公認解釋。
\end{延伸框}

\section*{\sffamily 5-6\quad 本章重點整理}

\begin{補充框}{一頁複習（公式與關鍵結論）}
\begin{itemize}\setlength{\itemsep}{2pt}
\item \textbf{能量守恆}：總能量不生不滅，只在各形式間轉換，$E_{\text{總}}=$ 定值。能量形式（動能／位能／熱／化學／電磁／核能）的劃分\textbf{隨尺度而變}（熱能微觀上就是分子動能）。電磁場帶能量，電磁波（光、訊號）能在真空中搬運能量。
\item \textbf{質能等價}：$E=mc^2$，$c^2=9\times10^{16}$ 極大，故微小質量對應巨大能量。質量本身是能量的一種形式；任何放能過程質量都會微減（質量虧損 $\Delta m$，放出能量 $\Delta m\,c^2$）。
\item \textbf{核能}：束縛能曲線在\textbf{鐵}達峰；\textbf{輕核融合（左）與重核分裂（右）都釋能}，皆往鐵峰靠攏。太陽靠融合，核電靠分裂。
\item \textbf{溫度／內能／熱}：溫度＝分子平均動能（理想氣體 $E_{\text{內}}\propto T$，須用克氏溫標 $T_K=t_C+273.15$）；內能＝物體內部所有微觀動能＋交互作用能（物體\textbf{擁有}的）；熱＝因\textbf{溫差}而\textbf{傳遞}中的能量（過程量，物體不「含有」熱）。熱只自發地由高溫流向低溫，溫度相等時達\textbf{熱平衡}；傳熱三途徑＝\textbf{傳導／對流／輻射}（只有輻射能穿越真空）。絕對零度 0 K 是溫度下限。
\item \textbf{第二定律（定性）}：功可 100\% 變熱，熱不能 100\% 變功；熱機效率 $\eta=1-\dfrac{Q_C}{Q_H}<1$。能量有\textbf{品質}，自然過程品質下降、熵增、不可逆，指向時間方向。
\end{itemize}
\end{補充框}

\end{document}
